\documentclass[12pt,a4paper]{article}

% ==================== 中文支持 ====================
\usepackage[UTF8]{ctex}

% ==================== 页面设置 ====================
\usepackage{geometry}
\geometry{left=2.5cm,right=2.5cm,top=2.5cm,bottom=2.5cm}

% ==================== 数学公式 ====================
\usepackage{amsmath,amssymb}

% ==================== 图表与浮动体 ====================
\usepackage{graphicx}
\usepackage{float}

% ==================== 表格 ====================
\usepackage{booktabs}
\usepackage{multirow}
\usepackage{array}

% ==================== 超链接 ====================
\usepackage{hyperref}
\hypersetup{
    colorlinks=true,
    linkcolor=blue,
    citecolor=blue,
    urlcolor=blue,
}

% ==================== 其他 ====================
\usepackage{enumitem}

\title{问题三：针对痰湿体质患者的个性化干预方案优化模型}
\author{参赛队伍}
\date{\today}

\begin{document}

\maketitle

\section{问题背景与建模思路}

\subsection{背景主旨：中西医结合的科学性}
高血脂症在中医理论体系中多归属于“痰湿”、“膏脂”范畴。其核心病机在于脾失健运，导致水湿内停，聚而成痰，这与现代医学中脂质代谢紊乱的病理机制高度契合。本模型旨在融合中医“治未病”的调理理念与西医运动康复的干预手段，通过量化“中医调理分级”（改善内环境）与“活动干预强度”（促进外代谢），构建针对中老年高血脂人群的精准、可行的个性化干预方案。

\subsection{核心建模逻辑}
针对问题三，我们构建了一个离散非线性优化模型，旨在有限的资源约束下，最大化干预效果。模型的核心要素定义如下：
\begin{itemize}
    \item \textbf{决策变量}：
    \begin{itemize}
        \item $x_1 \in \{1, 2, 3\}$：中医调理分级（1-基础，2-中度，3-强化），由患者初始痰湿积分决定，在静态模型中是参数，在动态模型中是阶段决策变量。
        \item $x_2 \in \{1, 2, 3\}$：活动干预强度（1-低，2-中，3-高）。
        \item $x_3 \in \{5,6,7,8,9,10\}$：每周活动干预频率（次/周），由于频率低于5次/周时疗效收益为零，而成本线性增长，因此每周5次是所有理性决策场景下的“保底频率”。
    \end{itemize}
    \item \textbf{目标函数}：在6个月的干预周期内，最大化痰湿积分的预期下降率。
    \item \textbf{约束条件}：严格遵循附件中关于年龄阶段、活动能力评分（ADL+IADL）对活动强度的限制，以及总经济成本不超过2000元的预算约束。
\end{itemize}

\section{静态基准方案模型建立与求解}

\subsection{疗效与成本量化公式}
根据附件附表2与附表3提供的信息，我们对干预效果与成本进行量化。

\paragraph{痰湿积分下降率 $R$ 的量化}
每月痰湿积分下降率 $R$ 与活动干预强度 $x_2$ 及频率 $x_3$ 的关系如下：
\begin{equation}
R =
\begin{cases}
 (x_2 - 1) \times 3\% + (x_3 - 5) \times 1\%, & x_3 \ge 5 \\
0, & x_3 < 5
\end{cases}
\end{equation}

经过6个月（24周）的干预后，预期的最终痰湿积分 $Score_{final}$ 计算公式为：
\begin{equation}
Score_{final} = Score_{initial} \times (1 - R)^6
\end{equation}

\paragraph{总成本计算公式}
6个月干预周期的总成本 $Cost_{total}$ 由两部分构成：固定调理成本和可变活动成本。
\begin{equation}
Cost_{total} = 6 \times Cost_{cond}(x_1) + 24 \times x_3 \times Cost_{act}(x_2)
\end{equation}
其中，$Cost_{cond}(x_1)$ 为表2中各级别调理的月成本，$Cost_{act}(x_2)$ 为表3中各级别活动的单次成本。

\subsection{模型求解方法}
由于本模型的决策变量均为离散变量，且组合空间有限（$3 \times 3 \times 10 = 90$ 种可能），我们采用穷举搜索法进行求解。该方法能够确保在严格满足所有约束条件的前提下，找到全局最优解。

\subsection{样本ID 1、2、3 求解结果}
依据上述模型，我们对附件中“样本ID”为1、2、3的三位确诊为痰湿体质的患者进行最优方案求解。通过编程遍历所有可行组合，得到如下基准结果：

\begin{table}[htbp]
\centering
\caption{样本ID 1、2、3的静态基准方案求解结果}
\label{tab:sample_results}
\begin{tabular}{cccccccc}
\toprule
\textbf{样本ID} & \textbf{最优调理} & \textbf{最优强度} & \textbf{最优频率} & \textbf{6个月总成本} & \textbf{最终积分} & \textbf{下降百分比} \\
& $x_1$ & $x_2$ & $x_3$ & (元) & (分) & \\
\midrule
1 & 3级 & 1级 & 10次 & 1500 & 47.05 & 26.49\% \\
2 & 1级 & 2级 & 10次 & 1380 & 35.17 & 39.36\% \\
3 & 2级 & 2级 & 10次 & 1680 & 35.77 & 39.36\% \\
\bottomrule
\end{tabular}
\end{table}

\paragraph{结果分析：}
    样本1的案例极具代表性。由于其活动能力评分低于40分，模型强制其活动强度 $x_2$ 锁定为1级。尽管患者采用了最高级别的强化调理（$x_1 = 3$），且活动频率达到上限（$x_3 = 10$次/周），其痰湿积分下降幅度依然受限。这深刻揭示了“体医融合”的基础性逻辑：个体的身体活动能力是实现最佳干预效果的必要前提。对于活动受限人群，干预的边际效益相对较低。

\section{模型扩展：多阶段动态调整分析（以样本1为例）}

\subsection{动态优化思想}
前述静态模型假设6个月内的干预方案一成不变。然而在实际干预过程中，随着患者痰湿积分的下降，其体质状态会发生改善。当积分降低并跨越附表2所定义的分级阈值时，对应的调理级别和成本理应进行动态调整。这种调整不仅更具经济性，也体现了精准医疗中“随证治之”的原则。我们以每2个月为一个阶段，对样本1进行动态路径推演。

\subsection{样本1的动态路径推演}
\begin{itemize}
    \item \textbf{阶段一（第1-2月）}：初始积分为64分，属于“强化调理（3级）”，月成本130元。在 $x_2=1, x_3=10$ 的最优方案下，每月下降率为5\%。
    \item \textbf{状态触发}：由于 $x_2=1, x_3=10$ 的下降率为每月 5\%。第2个月结束时，积分预计降至：
    \begin{equation}
    64 \times (1-0.05)^2 \approx 57.76 \text{分}
    \end{equation}
    \item \textbf{阶段二（第3-4月）}：此时积分57.76 < 59分，根据表2，调理分级可由3级下调至1级（基础调理）。月成本随之从130元降至30元。
    \item \textbf{阶段三（第5-6月）}：继续保持1级调理。
\end{itemize}

\subsection{动态效益}
\begin{itemize}
    \item \textbf{成本节约}：后4个月共节约 $(130-30) \times 4 = 400$ 元。
    \item \textbf{方案优化启示}：节约的成本可用于其他辅助性健康管理，或作为患者依从性的经济激励。更重要的是，该动态模型印证了中西医结合干预的“收敛性”与“精准性”。随着体质的改善，干预重心从“强化治疗”平滑过渡至“常态化健康管理”，有效避免了过度医疗，完美契合了中医“瘥后防复”的核心思想，即在病情好转后，通过低成本、可持续的调理方式防止复发。
\end{itemize}

\section{模型深化：基于多目标权衡的综合效用模型}

\subsection{综合效用函数的构建}
在基准模型中，2000元的预算上限对多数患者而言可能较为充裕。为探索“经济成本”与“健康疗效”之间的最优平衡，避免盲目追求指标改善，我们引入一个综合效用函数 $Z$ 进行多目标决策分析。

为解决成本（元）与疗效（积分下降量）量纲不同的问题，我们采用离差标准化方法对两者进行归一化处理，并构建如下最小化目标函数：
\begin{equation}
\min \; Z = \omega_1 \cdot \hat{C} - \omega_2 \cdot \hat{B}
\end{equation}
其中：
\begin{itemize}
    \item $\hat{C} = \frac{Cost_{total}}{Cost_{max}}$ 为归一化成本支出（$Cost_{max}=2000$）。
    \item $\hat{B} = \frac{\Delta Phlegm}{Phlegm_{initial}}$ 为归一化疗效收益，即痰湿积分下降百分比。
    \item $\omega_1, \omega_2$ 分别为成本和疗效的偏好权重系数，满足 $\omega_1 + \omega_2 = 1$，反映了决策者对“省钱”与“治病”的偏好权重，通过调整 $\omega$，可以模拟不同决策偏好（如医保控费导向或患者健康导向）下的最优方案。
\end{itemize}

\begin{table}[htbp]
\centering
\caption{不同偏好权重下的最优干预方案对比（以样本1为例）}
\label{tab:weight_preference}
\begin{tabular}{cccccl}
\toprule
\textbf{决策偏好} & $\boldsymbol{(\omega_1, \omega_2)}$ & $\boldsymbol{x_3}$ & \textbf{总成本(元)} & \textbf{下降率} & \textbf{决策逻辑} \\
\midrule
疗效优先 & (0.2, 0.8) & 10 & 1500 & 26.49\% & 牺牲经济性换取最大化康复速度 \\
平衡点   & (0.5, 0.5) & 10 & 1500 & 26.49\% & 1\%疗效的边际收益仍大于其边际成本 \\
成本优先 & (0.8, 0.2) & 5  & 1140 & 0.00\%  & 仅执行基础调理，停止高频活动干预 \\
\bottomrule
\end{tabular}
\end{table}

\subsection{模型灵敏度分析与讨论（以样本1为例）}
通过对样本1进行灵敏度分析，观察疗效权重 $\omega_2$ 在 $[0.1, 0.9]$ 区间内变化时最优干预频率 $x_3$ 与总成本的响应特征（如图1所示），可得出以下深刻结论：

\begin{enumerate}
    \item \textbf{低权重区的“稳健保守”策略}：当疗效权重 $\omega_2 \in [0.1, 0.38]$ 时，最优干预频率稳定在5次/周。这是因为每周5次是疗效产生的临界点。在此区间内，决策者（或模型）更看重经济成本，故倾向于选择能达到疗效的“最低有效投入”。这表明模型具有良好的鲁棒性，不会因微小偏好扰动而改变核心决策。
    
    \item \textbf{决策跃迁的“门槛效应”}：当 $\omega_2$ 跨越0.40附近时，最优频率出现显著“阶梯式”跃升，从5次/周迅速增至10次/周。此跃迁点标志着边际健康收益正式超过边际经济成本。它揭示了在实际慢病管理中，当患者的康复意愿或对健康的重视程度（即权重）超过某一阈值后，加大干预力度的方案便成为更优选择。
    
    \item \textbf{高权重区的“饱和干预”特征}：当 $\omega_2 > 0.45$ 后，最优频率锁定在10次/周，总成本也稳定在约1500元。此时决策已完全转为“疗效导向型”。尽管疗效权重继续增加，但由于受到每周10次的耐受度上限约束，干预强度已达模型允许的极限。
    
    \item \textbf{成本与频率的同步性}：图中总成本曲线与最优频率曲线呈现高度同步的变化趋势，这有力地证明了在调理等级固定的情况下，活动干预频率是驱动健康成本波动的核心变量。
\end{enumerate}

\begin{figure}[H]
\centering
\includegraphics[width=0.8\textwidth]{sensitivity_analysis.png}
\caption{疗效权重 $\omega_2$ 变化时最优干预频率与总成本的响应}
\label{fig:sensitivity}
\end{figure}

\section{模型完备形态与临床应用指南}

\subsection{模型集成与完备形态}
综合前述模型构建与深化分析，我们最终形成了一个多层级、可配置的决策支持系统框架。

\begin{figure}[H]  % 或使用 [htbp] 让 LaTeX 自动排版
\centering
\includegraphics[width=0.8\textwidth]{model.png}
\caption{模型完备形态的四层架构}
\label{fig:model_architecture}
\end{figure}

该完备模型具备以下核心能力：
\begin{enumerate}
    \item \textbf{个性化约束嵌入}：严格遵循患者年龄、活动能力等生理限制，确保方案安全可行。
    \item \textbf{动静结合优化}：既支持6个月固定方案的快速求解，亦支持基于积分变化的阶段性动态调整。
    \item \textbf{多目标决策支持}：通过综合效用函数，允许决策者在“经济成本”与“健康疗效”之间进行灵活权衡。
    \item \textbf{临床可解释性}：输出结果包含明确的特征-方案对应关系，便于基层医生理解与操作。
\end{enumerate}

\subsection{临床应用：“患者特征-最优方案”匹配规律表}
基于对附件数据中所有痰湿体质患者（体质标签=5）的模拟寻优结果，选取偏好权重系数均为0.5，并结合模型的普适性推演，我们归纳了不同特征组合下的最优干预方案推荐。下表旨在为基层医生提供直观、可查的临床决策参考。

\begin{center}
\includegraphics[width=\textwidth]{matching_table.png}  % 请替换为实际图片文件名

\vspace{0.5em}
\textbf{图 X \quad 痰湿体质患者个性化干预方案推荐表（偏好权重 $\omega_1=\omega_2=0.5$）}
\end{center}

{\small
\textbf{注：}
\begin{enumerate}[leftmargin=*, nosep]
    \item 活动能力总分为ADL与IADL评分之和（0-100分）。
    \item 6月后积分下降率基于每月4周、6个月共24周的干预模型估算，实际效果可能存在个体差异。
    \item 当活动能力<40分时，活动强度强制为1级，此时推荐通过最大频率（10次/周）补偿效果。
\end{enumerate}
}

\subsection{匹配规律的核心解读}

\paragraph{规律一：年龄与活动能力共同决定活动强度的上限。}
\begin{itemize}
    \item 40-59岁且活动能力≥60分者，可耐受最高3级强度，效果最佳。
    \item 60-79岁者，活动强度上限为2级，即便活动能力评分再高，亦不推荐3级强度。
    \item 80岁及以上或活动能力<40分者，活动强度严格限制为1级，安全优先。
\end{itemize}

\paragraph{规律二：活动受限时，高频次是唯一的补偿路径。}
\begin{itemize}
    \item 对于活动能力<40分的患者（如样本1），由于无法提升强度，模型无一例外地将周频率推至上限（10次/周）。这印证了“身体越受限，越需高频低强度运动”的干预原则。
\end{itemize}

\paragraph{规律三：初始痰湿积分决定调理起点，预算约束下调理等级可适时下调。}
\begin{itemize}
    \item 积分≥62分者必须从强化调理（3级）起步。
    \item 积分59-61分者适用中度调理（2级）。
    \item 积分≤58分者仅需基础调理（1级）。
    \item 当动态优化启动时，调理等级可随积分下降而降低，实现成本节约与精准治疗。
\end{itemize}

\paragraph{规律四：预算约束（2000元）在绝大多数情形下非紧约束。}
\begin{itemize}
    \item 即便在最高成本组合（3级调理+3级强度+10次/周）下，6个月总成本约为：
\begin{equation}
6 \times 130 + 24 \times 10 \times 8 = 780 + 1920 = 2700 \text{元}
\end{equation}
超出预算。因此对于40-59岁且活动能力≥60分的患者，若初始积分≥62分，模型将优先选择 3级强度+9次/周 或 2级强度+10次/周 的组合以符合预算要求。上表已根据预算约束进行了精细化调整。
\end{itemize}

\subsection{临床使用说明}
医生在使用本表时，可按以下步骤快速定位推荐方案：
\begin{enumerate}
    \item 确定年龄组：40-59岁、60-79岁或80-89岁。
    \item 计算活动能力总分：将患者的ADL和IADL评分相加。
    \item 查阅初始痰湿积分：根据中医体质评估结果获得。
    \item 在表中查找对应行，获取推荐调理等级、活动强度及周频率。
\end{enumerate}

例如，一位65岁男性患者，活动能力总分55分，初始痰湿积分63分。查表可得序号9推荐方案：3级调理+2级强度+9次/周，预期6个月积分下降约23\%。

本表为基于模型求解的标准化推荐，临床应用时仍需结合患者具体意愿、生活习惯及耐受度进行个体化微调。模型的动态版本（阶段优化）可进一步帮助医生在随访时根据积分变化灵活调整方案。

\section{结论与拓展分析}

\subsection{匹配规律总结}
\begin{itemize}
    \item \textbf{耐受度优先原则}：患者的基线活动能力（ADL/IADL评分）是决定干预方案上限的首要约束条件。对于活动能力受限人群，增加训练频率 $x_3$ 是弥补强度不足、提升干预效果的唯一可行路径。
    \item \textbf{经济边际效应评估}：在给定的2000元预算内，由于单位活动成本相对较低，模型倾向于在耐受度允许范围内选择较高频率。这同时也提示，对于预算紧张或依从性不高的患者，可适当降低频率，而不会对疗效产生颠覆性影响。
\end{itemize}

\subsection{中西医结合的深度反思}
本模型从数学量化的角度，充分证明了中西医结合干预的优越性：
\begin{enumerate}
    \item \textbf{协同增效机制}：中医调理（改善内环境）与西医活动干预（提升外代谢）之间存在显著的协同效应。两者结合的疗效曲线斜率，理论上优于任何单一干预方式。
    \item \textbf{模型未来拓展方向}：为使模型更具临床实用价值，未来可引入“疲劳度函数”或“依从性衰减模型”。例如，长期维持每周10次的高频训练可能导致中老年患者心理倦怠，从而降低实际依从性。引入基于心理学量表的动态依从性修正因子，将使模型的预测更为精准，方案推荐更贴合真实世界场景。
\end{enumerate}

\end{document}